% Auto-generated by build_annotation_scheme_table.py; do not edit by hand.
% Requires: booktabs, longtable, array
%
% % Auto-generated by build_annotation_scheme_table.py; do not edit by hand.
% Requires: booktabs, longtable, array
%
% % Auto-generated by build_annotation_scheme_table.py; do not edit by hand.
% Requires: booktabs, longtable, array
%
% % Auto-generated by build_annotation_scheme_table.py; do not edit by hand.
% Requires: booktabs, longtable, array
%
% \input{analysis/manuscript_tables/tex/annotation_scheme_Q1_Q2_Q3_Q5.tex}
%
% Reference examples: \ref{A}, \ref{B3}, \ref{D4a}

{\footnotesize
\setlength{\tabcolsep}{4pt}
\renewcommand{\arraystretch}{1.25}
\onecolumn\begin{longtable}{@{}>{\raggedright\arraybackslash}p{0.13\textwidth} >{\raggedright\arraybackslash}p{0.06\textwidth} >{\raggedright\arraybackslash}p{0.17\textwidth} >{\raggedright\arraybackslash}p{0.56\textwidth} c@{}}
        \toprule
        \textbf{Category} & \textbf{Code} & \textbf{Aspect} & \textbf{Definition} & \textbf{Pol.} \\
        \midrule
\endfirsthead
        \multicolumn{5}{c}{\textbf{Table \thetable{} -- continued from previous page}} \\
        \toprule
        \textbf{Category} & \textbf{Code} & \textbf{Aspect} & \textbf{Definition} & \textbf{Pol.} \\
        \midrule
\endhead
        \midrule
        \multicolumn{5}{r}{\textit{Continued on next page}} \\
        \midrule
\endfoot
        \bottomrule
        \noalign{\vskip 0.75em}
        \caption{Annotation scheme for participant comments on Single Q5, Q6, Compar. Q4, and Chunk Justif. Only positive codes are applied to comments for Single Q5, and only negative codes are applied for Single Q6. For Compar. Q4 and Chunk Justif., positive codes are applied to match the comments' positive aspects for the preferred translation, and negative codes applied to match the negative aspects of the non-preferred translation.}
        \label{tab:annotation-scheme-q1-q2-q3-q5} \\
\endlastfoot
        \parbox[t]{0.13\textwidth}{\codelabel{A}.\\Language-level features} & \codelabel{A1} & Grammar \& spelling &
        Issues related to grammar such as verb-subject agreement, article use, or misspellings. &
        $+/-$ \\[5pt]
         & \codelabel{A1b} & Tense &
        Tense correctness and consistency, including complex tense forms (past perfect for sequencing, sequence-of-tenses) and discourse-level tense/time-frame switching across the text. &
        $+/-$ \\[5pt]
         & \codelabel{A2} & Word choice: clarity &
        Whether individual word and phrase choices are clear or keep the reader guessing when it was not inteded. &
        $+/-$ \\[5pt]
         & \codelabel{A3} & Word choice: naturalness &
        Whether the choice of specific words or phrases is natural in English at the local level. Covers words/phrases that are grammatically correct and even undertandable, but sound off (e.g., awkward) in English. &
        $+/-$ \\[5pt]
         & \codelabel{A4} & Word choice: idiomaticity &
        Correct use of English idioms, set phrases, conventional collocations, and figure-of-speech precision; the negative side covers both misused/garbled idioms and over-reliance on clich\'{e}d or trite expressions. &
        $+/-$ \\[5pt]
         & \codelabel{A5} & Word choice: register \& contextual fit &
        Whether word choices match the context in aspects such as formality, period (e.g., historical vs contemporary), dialect/local speech, genre (e.g., sci-fi), speaker/situation (e.g., emotional context). &
        $+/-$ \\[5pt]
         & \codelabel{A6} & Word choice: richness vs. blandness &
        Lexical interest of the vocabulary used, such as unusual, varied, atmospheric words and turns of phrase vs.\ generic, common, sterile, ``dumbed-down'' choices. Independent of B9 (a passage can be rich-and-concise or rich-and-wordy). &
        $+/-$ \\[5pt]
         & \codelabel{A7} & Cultural adaptation &
        Translator's localization decisions in both directions: localizing/explicating foreign concepts for the target reader (e.g., glossing terms, Americanizing place names) vs.\ preserving foreignness (leaving names/terms in source language, retaining cultural markers). &
        $+/-$ \\[5pt]
         & \codelabel{A8} & Sentence structure &
        Sentence length, complexity, word order, run-ons, fragments, choppiness, and variety vs.\ monotony across sentences. &
        $+/-$ \\[5pt]
         & \codelabel{A9} & Cadence / rhythm / tempo &
        Prose rythm (whether the sentences have flow that feels natural and pleasible), tempo that is matching scenes content, etc. &
        $+/-$ \\[5pt]
         & \codelabel{A10} & Consistency &
        Whether the translation remains the same choices across the text in terms of vocabulary, voice, register, etc.\ (e.g., whether character names are translated consistently). &
        $+/-$ \\[5pt]
         & \codelabel{A11} & Formatting \& typography &
        The text's visual and typographic conventions, such as the typesetter-and-stylistic-choice layer that sits between grammar and content. It includes paragraph structure (breaks, dividers), dialogue formatting (quotation vs em-dashes). &
        $+/-$ \\[5pt]
         & \codelabel{A12} & Repetition or redundancy &
        Overuse of the same words, phrases, structures whether variation or compression would be better. This includes unnecessary lexical repetition, redundant content, etc. &
        $-$ \\[5pt]
         & \codelabel{A13} & Untranslated content &
        Source material (words or phrases) that appear untranslated in the target text and are perceived by the reader as translation failures not deliberate choices to retain the source word for culture preservation. &
        $-$ \\[5pt]
        \midrule
        \parbox[t]{0.13\textwidth}{\codelabel{B}.\\Narrative-level features} & \codelabel{B1} & Dialogue: naturalness &
        Whether the dialogue sounds like how people actually speak in the specific situation, context, emotional moment. It covers cases such as word plausability in speech (of this character), belivable phrasing etc. &
        $+/-$ \\[5pt]
         & \codelabel{B2} & Character voice &
        Whether characters come through as distinct, identifiable, and well-matched to who they are. B2 evaluates the voice of characters and narrators: both in dialogue and in prose attributed to a character's perspective. character distinctness, attribution \& character fit &
        $+/-$ \\[5pt]
         & \codelabel{B3} & Character portrayal &
        Whether characters feel fleshed out, developed, and relatable: depth of inner life and personality as expressed through the prose, beyond just voice. &
        $+/-$ \\[5pt]
         & \codelabel{B4} & Description / vividness / visualizability &
        Whether the reader can picture the scene and whether word choices conjure clear mental images even when both options are technically correct; includes the craft distinction of showing vs.\ telling (scene-rendering that creates imagery vs.\ bare informational delivery). &
        $+/-$ \\[5pt]
         & \codelabel{B5} & Figurative language &
        Use and effectiveness of metaphors, similes, imagery, and atmospheric or evocative phrasings; fires positively for evocative figures and negatively for misfiring ones that don't match the referent. &
        $+/-$ \\[5pt]
         & \codelabel{B6} & Emotional conveyance &
        Effectiveness of conveying the character’s or scene’s emotions to the reader (distinct from C5: B6 is about how well the reader can feel the story’s emotional elements, while C5 is about whether the text itself feels alive or soulful). &
        $+/-$ \\[5pt]
         & \codelabel{B7} & Narrative organization &
        Story structure: scene transitions, paragraph / section flow, balance of exposition vs.\ action, story cohesion / motif recall (themes and descriptions that call back to previous text in the story). &
        $+/-$ \\[5pt]
         & \codelabel{B8} & Narration \& POV effects &
        Narrator stance: interior monologue, free indirect discourse (blending character’s thoughts into narration), narrator’s distance from the action, narrator addressing the reader directly; point-of-view handling like POV shifts. &
        $+/-$ \\[5pt]
         & \codelabel{B9} & Pacing / information delivery / wordiness &
        Narrative-level pacing: info-dumps, lengthy descriptions that interrupt action; phrase-level wordiness (e.g.\ one word vs.\ whole phrase to describe the same thing). Independent of A6. &
        $+/-$ \\[5pt]
        \midrule
        \parbox[t]{0.13\textwidth}{\codelabel{C}.\\Reader experience} & \codelabel{C1} & Comprehension / ease of following &
        Whether the reader could understand what was happening: local sentence-level understanding and global scene/plot-level comprehension (these can fire independently). &
        $+/-$ \\[5pt]
         & \codelabel{C2} & Smoothness / reading effort &
        Whether the prose flowed or the reader kept getting stuck; i.e.\ how much friction was in the reading experience, separate from whether they understood it. &
        $+/-$ \\[5pt]
         & \codelabel{C3} & Engagement / immersion &
        Whether the reader was pulled into the story: being in the world, losing awareness of reading, focused attention on the scene. &
        $+/-$ \\[5pt]
         & \codelabel{C4} & Hook / pull-through &
        Whether the reader wanted to keep reading, find out what happens next, or continue past the given excerpt. &
        $+/-$ \\[5pt]
         & \codelabel{C5} & Humanness / soul &
        Holistic judgement on whether the prose feels alive and human vs.\ robotic or mechanical (distinct from A3: individual word choices or sentences can feel natural, while the whole story can overall feel soulless). &
        $+/-$ \\[5pt]
         & \codelabel{C6} & Enjoyment / overall positive affect &
        General reader satisfaction with the text: ``I liked it'', ``pleasant to read'', ``a pleasure'', ``I enjoyed both''; broader affective stance not related to a specific dimension, can be used as a catch-all for positive sentiment when no specific details given. &
        $+/-$ \\[5pt]
        \midrule
        \parbox[t]{0.13\textwidth}{\codelabel{D}.\\Meta-translation} & \codelabel{D1} & Doesn't read as translated &
        Holistic judgment on whether the prose reads as if originally written in English vs.\ feels like it’s been translated. &
        $+/-$ \\[5pt]
         & \codelabel{D2} & Adaptation vs. literalness &
        Whether the translator adapted to English idioms and general audience expectations (positive) vs.\ translated word-for-word from the source language (negative); can fire with D4b when literalness is interpreted as a cue for MT. &
        $+/-$ \\[5pt]
         & \codelabel{D3} & Faithfulness to original &
        How well the text preserves the source meaning, voice, emotion, structure, subtle implications, connotations, and specific facts; usually perceived at the chunk-level, and often inferred by comparing small detail inconsistencies between two presented translations. &
        $+/-$ \\[5pt]
         & \codelabel{D4a} & MT/AI verdict &
        Holistic reader feeling that the translation feels machine-generated or human, but without naming specific cues. &
        $+/-$ \\[5pt]
         & \codelabel{D4b} & Specific MT/AI tell named &
        Reader explicitly articulates the cue that triggers the MT/AI impression: e.g., em-dash overuse, bullet-point styling, ``would'' overuse, US English defaults, anthropomorphizing concepts, and other details that are percieved as tells for AI-generated text. &
        $-$ \\[5pt]
\end{longtable}
\twocolumn}

%
% Reference examples: \ref{A}, \ref{B3}, \ref{D4a}

{\footnotesize
\setlength{\tabcolsep}{4pt}
\renewcommand{\arraystretch}{1.25}
\onecolumn\begin{longtable}{@{}>{\raggedright\arraybackslash}p{0.13\textwidth} >{\raggedright\arraybackslash}p{0.06\textwidth} >{\raggedright\arraybackslash}p{0.17\textwidth} >{\raggedright\arraybackslash}p{0.56\textwidth} c@{}}
        \toprule
        \textbf{Category} & \textbf{Code} & \textbf{Aspect} & \textbf{Definition} & \textbf{Pol.} \\
        \midrule
\endfirsthead
        \multicolumn{5}{c}{\textbf{Table \thetable{} -- continued from previous page}} \\
        \toprule
        \textbf{Category} & \textbf{Code} & \textbf{Aspect} & \textbf{Definition} & \textbf{Pol.} \\
        \midrule
\endhead
        \midrule
        \multicolumn{5}{r}{\textit{Continued on next page}} \\
        \midrule
\endfoot
        \bottomrule
        \noalign{\vskip 0.75em}
        \caption{Annotation scheme for participant comments on Single Q5, Q6, Compar. Q4, and Chunk Justif. Only positive codes are applied to comments for Single Q5, and only negative codes are applied for Single Q6. For Compar. Q4 and Chunk Justif., positive codes are applied to match the comments' positive aspects for the preferred translation, and negative codes applied to match the negative aspects of the non-preferred translation.}
        \label{tab:annotation-scheme-q1-q2-q3-q5} \\
\endlastfoot
        \parbox[t]{0.13\textwidth}{\codelabel{A}.\\Language-level features} & \codelabel{A1} & Grammar \& spelling &
        Issues related to grammar such as verb-subject agreement, article use, or misspellings. &
        $+/-$ \\[5pt]
         & \codelabel{A1b} & Tense &
        Tense correctness and consistency, including complex tense forms (past perfect for sequencing, sequence-of-tenses) and discourse-level tense/time-frame switching across the text. &
        $+/-$ \\[5pt]
         & \codelabel{A2} & Word choice: clarity &
        Whether individual word and phrase choices are clear or keep the reader guessing when it was not inteded. &
        $+/-$ \\[5pt]
         & \codelabel{A3} & Word choice: naturalness &
        Whether the choice of specific words or phrases is natural in English at the local level. Covers words/phrases that are grammatically correct and even undertandable, but sound off (e.g., awkward) in English. &
        $+/-$ \\[5pt]
         & \codelabel{A4} & Word choice: idiomaticity &
        Correct use of English idioms, set phrases, conventional collocations, and figure-of-speech precision; the negative side covers both misused/garbled idioms and over-reliance on clich\'{e}d or trite expressions. &
        $+/-$ \\[5pt]
         & \codelabel{A5} & Word choice: register \& contextual fit &
        Whether word choices match the context in aspects such as formality, period (e.g., historical vs contemporary), dialect/local speech, genre (e.g., sci-fi), speaker/situation (e.g., emotional context). &
        $+/-$ \\[5pt]
         & \codelabel{A6} & Word choice: richness vs. blandness &
        Lexical interest of the vocabulary used, such as unusual, varied, atmospheric words and turns of phrase vs.\ generic, common, sterile, ``dumbed-down'' choices. Independent of B9 (a passage can be rich-and-concise or rich-and-wordy). &
        $+/-$ \\[5pt]
         & \codelabel{A7} & Cultural adaptation &
        Translator's localization decisions in both directions: localizing/explicating foreign concepts for the target reader (e.g., glossing terms, Americanizing place names) vs.\ preserving foreignness (leaving names/terms in source language, retaining cultural markers). &
        $+/-$ \\[5pt]
         & \codelabel{A8} & Sentence structure &
        Sentence length, complexity, word order, run-ons, fragments, choppiness, and variety vs.\ monotony across sentences. &
        $+/-$ \\[5pt]
         & \codelabel{A9} & Cadence / rhythm / tempo &
        Prose rythm (whether the sentences have flow that feels natural and pleasible), tempo that is matching scenes content, etc. &
        $+/-$ \\[5pt]
         & \codelabel{A10} & Consistency &
        Whether the translation remains the same choices across the text in terms of vocabulary, voice, register, etc.\ (e.g., whether character names are translated consistently). &
        $+/-$ \\[5pt]
         & \codelabel{A11} & Formatting \& typography &
        The text's visual and typographic conventions, such as the typesetter-and-stylistic-choice layer that sits between grammar and content. It includes paragraph structure (breaks, dividers), dialogue formatting (quotation vs em-dashes). &
        $+/-$ \\[5pt]
         & \codelabel{A12} & Repetition or redundancy &
        Overuse of the same words, phrases, structures whether variation or compression would be better. This includes unnecessary lexical repetition, redundant content, etc. &
        $-$ \\[5pt]
         & \codelabel{A13} & Untranslated content &
        Source material (words or phrases) that appear untranslated in the target text and are perceived by the reader as translation failures not deliberate choices to retain the source word for culture preservation. &
        $-$ \\[5pt]
        \midrule
        \parbox[t]{0.13\textwidth}{\codelabel{B}.\\Narrative-level features} & \codelabel{B1} & Dialogue: naturalness &
        Whether the dialogue sounds like how people actually speak in the specific situation, context, emotional moment. It covers cases such as word plausability in speech (of this character), belivable phrasing etc. &
        $+/-$ \\[5pt]
         & \codelabel{B2} & Character voice &
        Whether characters come through as distinct, identifiable, and well-matched to who they are. B2 evaluates the voice of characters and narrators: both in dialogue and in prose attributed to a character's perspective. character distinctness, attribution \& character fit &
        $+/-$ \\[5pt]
         & \codelabel{B3} & Character portrayal &
        Whether characters feel fleshed out, developed, and relatable: depth of inner life and personality as expressed through the prose, beyond just voice. &
        $+/-$ \\[5pt]
         & \codelabel{B4} & Description / vividness / visualizability &
        Whether the reader can picture the scene and whether word choices conjure clear mental images even when both options are technically correct; includes the craft distinction of showing vs.\ telling (scene-rendering that creates imagery vs.\ bare informational delivery). &
        $+/-$ \\[5pt]
         & \codelabel{B5} & Figurative language &
        Use and effectiveness of metaphors, similes, imagery, and atmospheric or evocative phrasings; fires positively for evocative figures and negatively for misfiring ones that don't match the referent. &
        $+/-$ \\[5pt]
         & \codelabel{B6} & Emotional conveyance &
        Effectiveness of conveying the character’s or scene’s emotions to the reader (distinct from C5: B6 is about how well the reader can feel the story’s emotional elements, while C5 is about whether the text itself feels alive or soulful). &
        $+/-$ \\[5pt]
         & \codelabel{B7} & Narrative organization &
        Story structure: scene transitions, paragraph / section flow, balance of exposition vs.\ action, story cohesion / motif recall (themes and descriptions that call back to previous text in the story). &
        $+/-$ \\[5pt]
         & \codelabel{B8} & Narration \& POV effects &
        Narrator stance: interior monologue, free indirect discourse (blending character’s thoughts into narration), narrator’s distance from the action, narrator addressing the reader directly; point-of-view handling like POV shifts. &
        $+/-$ \\[5pt]
         & \codelabel{B9} & Pacing / information delivery / wordiness &
        Narrative-level pacing: info-dumps, lengthy descriptions that interrupt action; phrase-level wordiness (e.g.\ one word vs.\ whole phrase to describe the same thing). Independent of A6. &
        $+/-$ \\[5pt]
        \midrule
        \parbox[t]{0.13\textwidth}{\codelabel{C}.\\Reader experience} & \codelabel{C1} & Comprehension / ease of following &
        Whether the reader could understand what was happening: local sentence-level understanding and global scene/plot-level comprehension (these can fire independently). &
        $+/-$ \\[5pt]
         & \codelabel{C2} & Smoothness / reading effort &
        Whether the prose flowed or the reader kept getting stuck; i.e.\ how much friction was in the reading experience, separate from whether they understood it. &
        $+/-$ \\[5pt]
         & \codelabel{C3} & Engagement / immersion &
        Whether the reader was pulled into the story: being in the world, losing awareness of reading, focused attention on the scene. &
        $+/-$ \\[5pt]
         & \codelabel{C4} & Hook / pull-through &
        Whether the reader wanted to keep reading, find out what happens next, or continue past the given excerpt. &
        $+/-$ \\[5pt]
         & \codelabel{C5} & Humanness / soul &
        Holistic judgement on whether the prose feels alive and human vs.\ robotic or mechanical (distinct from A3: individual word choices or sentences can feel natural, while the whole story can overall feel soulless). &
        $+/-$ \\[5pt]
         & \codelabel{C6} & Enjoyment / overall positive affect &
        General reader satisfaction with the text: ``I liked it'', ``pleasant to read'', ``a pleasure'', ``I enjoyed both''; broader affective stance not related to a specific dimension, can be used as a catch-all for positive sentiment when no specific details given. &
        $+/-$ \\[5pt]
        \midrule
        \parbox[t]{0.13\textwidth}{\codelabel{D}.\\Meta-translation} & \codelabel{D1} & Doesn't read as translated &
        Holistic judgment on whether the prose reads as if originally written in English vs.\ feels like it’s been translated. &
        $+/-$ \\[5pt]
         & \codelabel{D2} & Adaptation vs. literalness &
        Whether the translator adapted to English idioms and general audience expectations (positive) vs.\ translated word-for-word from the source language (negative); can fire with D4b when literalness is interpreted as a cue for MT. &
        $+/-$ \\[5pt]
         & \codelabel{D3} & Faithfulness to original &
        How well the text preserves the source meaning, voice, emotion, structure, subtle implications, connotations, and specific facts; usually perceived at the chunk-level, and often inferred by comparing small detail inconsistencies between two presented translations. &
        $+/-$ \\[5pt]
         & \codelabel{D4a} & MT/AI verdict &
        Holistic reader feeling that the translation feels machine-generated or human, but without naming specific cues. &
        $+/-$ \\[5pt]
         & \codelabel{D4b} & Specific MT/AI tell named &
        Reader explicitly articulates the cue that triggers the MT/AI impression: e.g., em-dash overuse, bullet-point styling, ``would'' overuse, US English defaults, anthropomorphizing concepts, and other details that are percieved as tells for AI-generated text. &
        $-$ \\[5pt]
\end{longtable}
\twocolumn}

%
% Reference examples: \ref{A}, \ref{B3}, \ref{D4a}

{\footnotesize
\setlength{\tabcolsep}{4pt}
\renewcommand{\arraystretch}{1.25}
\onecolumn\begin{longtable}{@{}>{\raggedright\arraybackslash}p{0.13\textwidth} >{\raggedright\arraybackslash}p{0.06\textwidth} >{\raggedright\arraybackslash}p{0.17\textwidth} >{\raggedright\arraybackslash}p{0.56\textwidth} c@{}}
        \toprule
        \textbf{Category} & \textbf{Code} & \textbf{Aspect} & \textbf{Definition} & \textbf{Pol.} \\
        \midrule
\endfirsthead
        \multicolumn{5}{c}{\textbf{Table \thetable{} -- continued from previous page}} \\
        \toprule
        \textbf{Category} & \textbf{Code} & \textbf{Aspect} & \textbf{Definition} & \textbf{Pol.} \\
        \midrule
\endhead
        \midrule
        \multicolumn{5}{r}{\textit{Continued on next page}} \\
        \midrule
\endfoot
        \bottomrule
        \noalign{\vskip 0.75em}
        \caption{Annotation scheme for participant comments on Single Q5, Q6, Compar. Q4, and Chunk Justif. Only positive codes are applied to comments for Single Q5, and only negative codes are applied for Single Q6. For Compar. Q4 and Chunk Justif., positive codes are applied to match the comments' positive aspects for the preferred translation, and negative codes applied to match the negative aspects of the non-preferred translation.}
        \label{tab:annotation-scheme-q1-q2-q3-q5} \\
\endlastfoot
        \parbox[t]{0.13\textwidth}{\codelabel{A}.\\Language-level features} & \codelabel{A1} & Grammar \& spelling &
        Issues related to grammar such as verb-subject agreement, article use, or misspellings. &
        $+/-$ \\[5pt]
         & \codelabel{A1b} & Tense &
        Tense correctness and consistency, including complex tense forms (past perfect for sequencing, sequence-of-tenses) and discourse-level tense/time-frame switching across the text. &
        $+/-$ \\[5pt]
         & \codelabel{A2} & Word choice: clarity &
        Whether individual word and phrase choices are clear or keep the reader guessing when it was not inteded. &
        $+/-$ \\[5pt]
         & \codelabel{A3} & Word choice: naturalness &
        Whether the choice of specific words or phrases is natural in English at the local level. Covers words/phrases that are grammatically correct and even undertandable, but sound off (e.g., awkward) in English. &
        $+/-$ \\[5pt]
         & \codelabel{A4} & Word choice: idiomaticity &
        Correct use of English idioms, set phrases, conventional collocations, and figure-of-speech precision; the negative side covers both misused/garbled idioms and over-reliance on clich\'{e}d or trite expressions. &
        $+/-$ \\[5pt]
         & \codelabel{A5} & Word choice: register \& contextual fit &
        Whether word choices match the context in aspects such as formality, period (e.g., historical vs contemporary), dialect/local speech, genre (e.g., sci-fi), speaker/situation (e.g., emotional context). &
        $+/-$ \\[5pt]
         & \codelabel{A6} & Word choice: richness vs. blandness &
        Lexical interest of the vocabulary used, such as unusual, varied, atmospheric words and turns of phrase vs.\ generic, common, sterile, ``dumbed-down'' choices. Independent of B9 (a passage can be rich-and-concise or rich-and-wordy). &
        $+/-$ \\[5pt]
         & \codelabel{A7} & Cultural adaptation &
        Translator's localization decisions in both directions: localizing/explicating foreign concepts for the target reader (e.g., glossing terms, Americanizing place names) vs.\ preserving foreignness (leaving names/terms in source language, retaining cultural markers). &
        $+/-$ \\[5pt]
         & \codelabel{A8} & Sentence structure &
        Sentence length, complexity, word order, run-ons, fragments, choppiness, and variety vs.\ monotony across sentences. &
        $+/-$ \\[5pt]
         & \codelabel{A9} & Cadence / rhythm / tempo &
        Prose rythm (whether the sentences have flow that feels natural and pleasible), tempo that is matching scenes content, etc. &
        $+/-$ \\[5pt]
         & \codelabel{A10} & Consistency &
        Whether the translation remains the same choices across the text in terms of vocabulary, voice, register, etc.\ (e.g., whether character names are translated consistently). &
        $+/-$ \\[5pt]
         & \codelabel{A11} & Formatting \& typography &
        The text's visual and typographic conventions, such as the typesetter-and-stylistic-choice layer that sits between grammar and content. It includes paragraph structure (breaks, dividers), dialogue formatting (quotation vs em-dashes). &
        $+/-$ \\[5pt]
         & \codelabel{A12} & Repetition or redundancy &
        Overuse of the same words, phrases, structures whether variation or compression would be better. This includes unnecessary lexical repetition, redundant content, etc. &
        $-$ \\[5pt]
         & \codelabel{A13} & Untranslated content &
        Source material (words or phrases) that appear untranslated in the target text and are perceived by the reader as translation failures not deliberate choices to retain the source word for culture preservation. &
        $-$ \\[5pt]
        \midrule
        \parbox[t]{0.13\textwidth}{\codelabel{B}.\\Narrative-level features} & \codelabel{B1} & Dialogue: naturalness &
        Whether the dialogue sounds like how people actually speak in the specific situation, context, emotional moment. It covers cases such as word plausability in speech (of this character), belivable phrasing etc. &
        $+/-$ \\[5pt]
         & \codelabel{B2} & Character voice &
        Whether characters come through as distinct, identifiable, and well-matched to who they are. B2 evaluates the voice of characters and narrators: both in dialogue and in prose attributed to a character's perspective. character distinctness, attribution \& character fit &
        $+/-$ \\[5pt]
         & \codelabel{B3} & Character portrayal &
        Whether characters feel fleshed out, developed, and relatable: depth of inner life and personality as expressed through the prose, beyond just voice. &
        $+/-$ \\[5pt]
         & \codelabel{B4} & Description / vividness / visualizability &
        Whether the reader can picture the scene and whether word choices conjure clear mental images even when both options are technically correct; includes the craft distinction of showing vs.\ telling (scene-rendering that creates imagery vs.\ bare informational delivery). &
        $+/-$ \\[5pt]
         & \codelabel{B5} & Figurative language &
        Use and effectiveness of metaphors, similes, imagery, and atmospheric or evocative phrasings; fires positively for evocative figures and negatively for misfiring ones that don't match the referent. &
        $+/-$ \\[5pt]
         & \codelabel{B6} & Emotional conveyance &
        Effectiveness of conveying the character’s or scene’s emotions to the reader (distinct from C5: B6 is about how well the reader can feel the story’s emotional elements, while C5 is about whether the text itself feels alive or soulful). &
        $+/-$ \\[5pt]
         & \codelabel{B7} & Narrative organization &
        Story structure: scene transitions, paragraph / section flow, balance of exposition vs.\ action, story cohesion / motif recall (themes and descriptions that call back to previous text in the story). &
        $+/-$ \\[5pt]
         & \codelabel{B8} & Narration \& POV effects &
        Narrator stance: interior monologue, free indirect discourse (blending character’s thoughts into narration), narrator’s distance from the action, narrator addressing the reader directly; point-of-view handling like POV shifts. &
        $+/-$ \\[5pt]
         & \codelabel{B9} & Pacing / information delivery / wordiness &
        Narrative-level pacing: info-dumps, lengthy descriptions that interrupt action; phrase-level wordiness (e.g.\ one word vs.\ whole phrase to describe the same thing). Independent of A6. &
        $+/-$ \\[5pt]
        \midrule
        \parbox[t]{0.13\textwidth}{\codelabel{C}.\\Reader experience} & \codelabel{C1} & Comprehension / ease of following &
        Whether the reader could understand what was happening: local sentence-level understanding and global scene/plot-level comprehension (these can fire independently). &
        $+/-$ \\[5pt]
         & \codelabel{C2} & Smoothness / reading effort &
        Whether the prose flowed or the reader kept getting stuck; i.e.\ how much friction was in the reading experience, separate from whether they understood it. &
        $+/-$ \\[5pt]
         & \codelabel{C3} & Engagement / immersion &
        Whether the reader was pulled into the story: being in the world, losing awareness of reading, focused attention on the scene. &
        $+/-$ \\[5pt]
         & \codelabel{C4} & Hook / pull-through &
        Whether the reader wanted to keep reading, find out what happens next, or continue past the given excerpt. &
        $+/-$ \\[5pt]
         & \codelabel{C5} & Humanness / soul &
        Holistic judgement on whether the prose feels alive and human vs.\ robotic or mechanical (distinct from A3: individual word choices or sentences can feel natural, while the whole story can overall feel soulless). &
        $+/-$ \\[5pt]
         & \codelabel{C6} & Enjoyment / overall positive affect &
        General reader satisfaction with the text: ``I liked it'', ``pleasant to read'', ``a pleasure'', ``I enjoyed both''; broader affective stance not related to a specific dimension, can be used as a catch-all for positive sentiment when no specific details given. &
        $+/-$ \\[5pt]
        \midrule
        \parbox[t]{0.13\textwidth}{\codelabel{D}.\\Meta-translation} & \codelabel{D1} & Doesn't read as translated &
        Holistic judgment on whether the prose reads as if originally written in English vs.\ feels like it’s been translated. &
        $+/-$ \\[5pt]
         & \codelabel{D2} & Adaptation vs. literalness &
        Whether the translator adapted to English idioms and general audience expectations (positive) vs.\ translated word-for-word from the source language (negative); can fire with D4b when literalness is interpreted as a cue for MT. &
        $+/-$ \\[5pt]
         & \codelabel{D3} & Faithfulness to original &
        How well the text preserves the source meaning, voice, emotion, structure, subtle implications, connotations, and specific facts; usually perceived at the chunk-level, and often inferred by comparing small detail inconsistencies between two presented translations. &
        $+/-$ \\[5pt]
         & \codelabel{D4a} & MT/AI verdict &
        Holistic reader feeling that the translation feels machine-generated or human, but without naming specific cues. &
        $+/-$ \\[5pt]
         & \codelabel{D4b} & Specific MT/AI tell named &
        Reader explicitly articulates the cue that triggers the MT/AI impression: e.g., em-dash overuse, bullet-point styling, ``would'' overuse, US English defaults, anthropomorphizing concepts, and other details that are percieved as tells for AI-generated text. &
        $-$ \\[5pt]
\end{longtable}
\twocolumn}

%
% Reference examples: \ref{A}, \ref{B3}, \ref{D4a}

{\footnotesize
\setlength{\tabcolsep}{4pt}
\renewcommand{\arraystretch}{1.25}
\onecolumn\begin{longtable}{@{}>{\raggedright\arraybackslash}p{0.13\textwidth} >{\raggedright\arraybackslash}p{0.06\textwidth} >{\raggedright\arraybackslash}p{0.17\textwidth} >{\raggedright\arraybackslash}p{0.56\textwidth} c@{}}
        \toprule
        \textbf{Category} & \textbf{Code} & \textbf{Aspect} & \textbf{Definition} & \textbf{Pol.} \\
        \midrule
\endfirsthead
        \multicolumn{5}{c}{\textbf{Table \thetable{} -- continued from previous page}} \\
        \toprule
        \textbf{Category} & \textbf{Code} & \textbf{Aspect} & \textbf{Definition} & \textbf{Pol.} \\
        \midrule
\endhead
        \midrule
        \multicolumn{5}{r}{\textit{Continued on next page}} \\
        \midrule
\endfoot
        \bottomrule
        \noalign{\vskip 0.75em}
        \caption{Annotation scheme for participant comments on Single Q5, Q6, Compar. Q4, and Chunk Justif. Only positive codes are applied to comments for Single Q5, and only negative codes are applied for Single Q6. For Compar. Q4 and Chunk Justif., positive codes are applied to match the comments' positive aspects for the preferred translation, and negative codes applied to match the negative aspects of the non-preferred translation.}
        \label{tab:annotation-scheme-q1-q2-q3-q5} \\
\endlastfoot
        \parbox[t]{0.13\textwidth}{\codelabel{A}.\\Language-level features} & \codelabel{A1} & Grammar \& spelling &
        Issues related to grammar such as verb-subject agreement, article use, or misspellings. &
        $+/-$ \\[5pt]
         & \codelabel{A1b} & Tense &
        Tense correctness and consistency, including complex tense forms (past perfect for sequencing, sequence-of-tenses) and discourse-level tense/time-frame switching across the text. &
        $+/-$ \\[5pt]
         & \codelabel{A2} & Word choice: clarity &
        Whether individual word and phrase choices are clear or keep the reader guessing when it was not inteded. &
        $+/-$ \\[5pt]
         & \codelabel{A3} & Word choice: naturalness &
        Whether the choice of specific words or phrases is natural in English at the local level. Covers words/phrases that are grammatically correct and even undertandable, but sound off (e.g., awkward) in English. &
        $+/-$ \\[5pt]
         & \codelabel{A4} & Word choice: idiomaticity &
        Correct use of English idioms, set phrases, conventional collocations, and figure-of-speech precision; the negative side covers both misused/garbled idioms and over-reliance on clich\'{e}d or trite expressions. &
        $+/-$ \\[5pt]
         & \codelabel{A5} & Word choice: register \& contextual fit &
        Whether word choices match the context in aspects such as formality, period (e.g., historical vs contemporary), dialect/local speech, genre (e.g., sci-fi), speaker/situation (e.g., emotional context). &
        $+/-$ \\[5pt]
         & \codelabel{A6} & Word choice: richness vs. blandness &
        Lexical interest of the vocabulary used, such as unusual, varied, atmospheric words and turns of phrase vs.\ generic, common, sterile, ``dumbed-down'' choices. Independent of B9 (a passage can be rich-and-concise or rich-and-wordy). &
        $+/-$ \\[5pt]
         & \codelabel{A7} & Cultural adaptation &
        Translator's localization decisions in both directions: localizing/explicating foreign concepts for the target reader (e.g., glossing terms, Americanizing place names) vs.\ preserving foreignness (leaving names/terms in source language, retaining cultural markers). &
        $+/-$ \\[5pt]
         & \codelabel{A8} & Sentence structure &
        Sentence length, complexity, word order, run-ons, fragments, choppiness, and variety vs.\ monotony across sentences. &
        $+/-$ \\[5pt]
         & \codelabel{A9} & Cadence / rhythm / tempo &
        Prose rythm (whether the sentences have flow that feels natural and pleasible), tempo that is matching scenes content, etc. &
        $+/-$ \\[5pt]
         & \codelabel{A10} & Consistency &
        Whether the translation remains the same choices across the text in terms of vocabulary, voice, register, etc.\ (e.g., whether character names are translated consistently). &
        $+/-$ \\[5pt]
         & \codelabel{A11} & Formatting \& typography &
        The text's visual and typographic conventions, such as the typesetter-and-stylistic-choice layer that sits between grammar and content. It includes paragraph structure (breaks, dividers), dialogue formatting (quotation vs em-dashes). &
        $+/-$ \\[5pt]
         & \codelabel{A12} & Repetition or redundancy &
        Overuse of the same words, phrases, structures whether variation or compression would be better. This includes unnecessary lexical repetition, redundant content, etc. &
        $-$ \\[5pt]
         & \codelabel{A13} & Untranslated content &
        Source material (words or phrases) that appear untranslated in the target text and are perceived by the reader as translation failures not deliberate choices to retain the source word for culture preservation. &
        $-$ \\[5pt]
        \midrule
        \parbox[t]{0.13\textwidth}{\codelabel{B}.\\Narrative-level features} & \codelabel{B1} & Dialogue: naturalness &
        Whether the dialogue sounds like how people actually speak in the specific situation, context, emotional moment. It covers cases such as word plausability in speech (of this character), belivable phrasing etc. &
        $+/-$ \\[5pt]
         & \codelabel{B2} & Character voice &
        Whether characters come through as distinct, identifiable, and well-matched to who they are. B2 evaluates the voice of characters and narrators: both in dialogue and in prose attributed to a character's perspective. character distinctness, attribution \& character fit &
        $+/-$ \\[5pt]
         & \codelabel{B3} & Character portrayal &
        Whether characters feel fleshed out, developed, and relatable: depth of inner life and personality as expressed through the prose, beyond just voice. &
        $+/-$ \\[5pt]
         & \codelabel{B4} & Description / vividness / visualizability &
        Whether the reader can picture the scene and whether word choices conjure clear mental images even when both options are technically correct; includes the craft distinction of showing vs.\ telling (scene-rendering that creates imagery vs.\ bare informational delivery). &
        $+/-$ \\[5pt]
         & \codelabel{B5} & Figurative language &
        Use and effectiveness of metaphors, similes, imagery, and atmospheric or evocative phrasings; fires positively for evocative figures and negatively for misfiring ones that don't match the referent. &
        $+/-$ \\[5pt]
         & \codelabel{B6} & Emotional conveyance &
        Effectiveness of conveying the character’s or scene’s emotions to the reader (distinct from C5: B6 is about how well the reader can feel the story’s emotional elements, while C5 is about whether the text itself feels alive or soulful). &
        $+/-$ \\[5pt]
         & \codelabel{B7} & Narrative organization &
        Story structure: scene transitions, paragraph / section flow, balance of exposition vs.\ action, story cohesion / motif recall (themes and descriptions that call back to previous text in the story). &
        $+/-$ \\[5pt]
         & \codelabel{B8} & Narration \& POV effects &
        Narrator stance: interior monologue, free indirect discourse (blending character’s thoughts into narration), narrator’s distance from the action, narrator addressing the reader directly; point-of-view handling like POV shifts. &
        $+/-$ \\[5pt]
         & \codelabel{B9} & Pacing / information delivery / wordiness &
        Narrative-level pacing: info-dumps, lengthy descriptions that interrupt action; phrase-level wordiness (e.g.\ one word vs.\ whole phrase to describe the same thing). Independent of A6. &
        $+/-$ \\[5pt]
        \midrule
        \parbox[t]{0.13\textwidth}{\codelabel{C}.\\Reader experience} & \codelabel{C1} & Comprehension / ease of following &
        Whether the reader could understand what was happening: local sentence-level understanding and global scene/plot-level comprehension (these can fire independently). &
        $+/-$ \\[5pt]
         & \codelabel{C2} & Smoothness / reading effort &
        Whether the prose flowed or the reader kept getting stuck; i.e.\ how much friction was in the reading experience, separate from whether they understood it. &
        $+/-$ \\[5pt]
         & \codelabel{C3} & Engagement / immersion &
        Whether the reader was pulled into the story: being in the world, losing awareness of reading, focused attention on the scene. &
        $+/-$ \\[5pt]
         & \codelabel{C4} & Hook / pull-through &
        Whether the reader wanted to keep reading, find out what happens next, or continue past the given excerpt. &
        $+/-$ \\[5pt]
         & \codelabel{C5} & Humanness / soul &
        Holistic judgement on whether the prose feels alive and human vs.\ robotic or mechanical (distinct from A3: individual word choices or sentences can feel natural, while the whole story can overall feel soulless). &
        $+/-$ \\[5pt]
         & \codelabel{C6} & Enjoyment / overall positive affect &
        General reader satisfaction with the text: ``I liked it'', ``pleasant to read'', ``a pleasure'', ``I enjoyed both''; broader affective stance not related to a specific dimension, can be used as a catch-all for positive sentiment when no specific details given. &
        $+/-$ \\[5pt]
        \midrule
        \parbox[t]{0.13\textwidth}{\codelabel{D}.\\Meta-translation} & \codelabel{D1} & Doesn't read as translated &
        Holistic judgment on whether the prose reads as if originally written in English vs.\ feels like it’s been translated. &
        $+/-$ \\[5pt]
         & \codelabel{D2} & Adaptation vs. literalness &
        Whether the translator adapted to English idioms and general audience expectations (positive) vs.\ translated word-for-word from the source language (negative); can fire with D4b when literalness is interpreted as a cue for MT. &
        $+/-$ \\[5pt]
         & \codelabel{D3} & Faithfulness to original &
        How well the text preserves the source meaning, voice, emotion, structure, subtle implications, connotations, and specific facts; usually perceived at the chunk-level, and often inferred by comparing small detail inconsistencies between two presented translations. &
        $+/-$ \\[5pt]
         & \codelabel{D4a} & MT/AI verdict &
        Holistic reader feeling that the translation feels machine-generated or human, but without naming specific cues. &
        $+/-$ \\[5pt]
         & \codelabel{D4b} & Specific MT/AI tell named &
        Reader explicitly articulates the cue that triggers the MT/AI impression: e.g., em-dash overuse, bullet-point styling, ``would'' overuse, US English defaults, anthropomorphizing concepts, and other details that are percieved as tells for AI-generated text. &
        $-$ \\[5pt]
\end{longtable}
\twocolumn}
